% Lösungen Einheit 3 von 4 – Antiproportionale Zuordnungen (zusammenbau v0.1)
\einheitenkopf{Einheit 3 von 4 · Antiproportionale Zuordnungen}

\erg{20}{a) weniger \quad b) 5 Tage \quad c) 3 h \quad d) 3 · 8 = 24 h für eine Maschine; 24 : 4 = 6 h \quad e) 6 · 10 = 60 = 4 · 15; es passt. \quad f) 3 · 4 = 12 Stunden für einen; 12 : 6 = 2 h \quad g) 6 · 10 = 60 Ziegentage; 60 : 5 = 12 Tage}
\erg{21}{a) 4 h; 6 Helfer}
\erg{22}{a) Kim hat proportional gerechnet; doppelt so viele Maler brauchen die halbe Zeit: 5 : 2 = 2,5 Tage. \quad b) Das Produkt ist die ganze Arbeit in Arbeitertagen; sie bleibt gleich, egal auf wie viele Arbeiter sie verteilt wird. \quad c) (1|6), (2|3), (3|2), (6|1), verbunden zu einer fallenden Kurve \quad d) 8 Helfer (5 · 3 = 15 Helferstunden, 15 : 2 = 7,5, aufrunden)}
